\documentclass[../main.tex]{subfiles}

\graphicspath{{../image/}{../figure/}}

\begin{document}

\chapter{采样}\label{chapter:ch14-采样}

在深度学习领域，我们经常需要从概率分布 $p(\bm{z})$ 中生成 $\bm{z}$ 的样本。这里的 $\bm{z}$ 可能是一个标量，分布可能是单变量 Gauss 分布；或者 $\bm{z}$ 可能是一幅高分辨率的图像，而 $p(\bm{z})$ 可能是由深度神经网络定义的生成式模型。对于许多简单的分布，我们可以使用数值技术直接生成合适的样本；而对于更复杂的分布，包括那些隐式定义的分布，我们则需要使用更复杂的方法。

\begin{definition}[采样与样本]\label{definition:ch14-采样}
设 $D\geqslant1$，概率空间 $(\Omega,\mathcal{F},\mathbb{P})$，随机变量 $\bm{Z}:\Omega\to\mathbb{R}^{D}$ 的分布为 $p(\bm{z})$（关于 Lebesgue 测度或计数测度）。若一个计算过程每次运行都返回 $\bm{Z}(\omega)$ 的一个实现值 $\bm{z}\in\mathbb{R}^{D}$，且对每个 Borel 集 $B\subseteq\mathbb{R}^{D}$ 有 $\mathbb{P}(\bm{Z}\in B)=\int_{B}p(\bm{z})\,\mathrm{d}\bm{z}$（离散情形为求和），则称该过程为从 $p$ \textbf{采样}（sampling），也称 Monte Carlo 采样（Monte Carlo sampling）；每次返回的实现值称为一个\textbf{样本}（sample）。
\end{definition}

% SUPPLEMENT: parent=definition:ch14-采样; kind=motivation; location=紧随其后的 remark
\begin{remark}本定义有两点约定值得说明。第一，本书把\emph{单个}被实例化的取值称为样本，而传统统计学常把一组观测称为一个样本；采用前一种约定后，$\bm{z}^{(l)}$ 这样的记号可以直接读作「第 $l$ 个样本」，本章的求和与平均才不会产生歧义。第二，定义只要求「生成取值」这一操作，并未要求 $p(\bm{z})$ 有解析表达式：这为本章后面处理未归一化分布 $p(\bm{z})=\widetilde{p}(\bm{z})/Z_p$（其中 $Z_p$ 未知）留下空间，因为从 $p$ 采样并不需要知道 $Z_p$。此外，计算机给出的样本实际上是伪随机数，它只保证通过统计检验意义上的随机性，一切基于样本的结论都带有统计误差，而\cref{subsection:ch14-期望}将用方差来度量这一误差。
\end{remark}

在本章中，我们将重点关注与深度学习最相关的采样方面的内容。关于 Monte Carlo 采样方法的更通用的介绍，可以参考 Gilks, Richardson, and Spiegelhalter (1996) 以及 Robert and Casella (1999)。

\subfile{Chapter-14/section-01.tex}
\subfile{Chapter-14/section-02.tex}
\subfile{Chapter-14/section-03.tex}
\subfile{Chapter-14/section-04.tex}

\end{document}
